% ============ 第2章 §2.6 Blaschke积与H^p的零点集(原书 pp.137-141) ============
\section{\texorpdfstring{Blaschke 积与 $H^p$ 的零点集}{Blaschke 积与 Hp 的零点集}}\label{sec:2.6}
\index{Blaschke 积!(单位圆板上的)}


当 $|a| < 1$ 且 $a \neq 0$ 时,$\frac{|a|}{a}\frac{a - z}{1 - \bar{a}z}$ 是从 $D$ 到 $D$ 上的一次变换,且
\begin{align*}
\left|1 - \frac{|a|}{a}\frac{a - z}{1 - \bar{a}z}\right| & = (1 - |a|)\frac{|a + |a|z|}{|a|\cdot|1 - \bar{a}z|} \\
& \leq (1 - |a|)\cdot\frac{2}{1 - |z|} \quad (z \in D).
\end{align*}
因此,设 $\{a_{j}\}_{j=1}^{\infty} \subset D \setminus \{0\}$,则
\begin{equation*}
\sum_{j = 1}^{\infty}\left|1 - \frac{|a_{j}|}{a_{j}}\frac{a_{j} - z}{1 - \bar{a}_{j}z}\right| \leq \frac{2}{1 - |z|}\sum_{j = 1}^{\infty}(1 - |a_{j}|) \quad (z \in D \setminus \{a_{j}\}_{j=1}^{\infty}),
\end{equation*}
所以
\begin{equation*}
\prod_{j = 1}^{\infty}\frac{|a_{j}|}{a_{j}}\frac{a_{j} - z}{1 - \bar{a}_{j}z}
\end{equation*}
对每个 $z \in D$ 绝对收敛的充要条件是
\begin{equation*}
\sum_{j = 1}^{\infty}(1 - |a_{j}|) < \infty
\end{equation*}
(必要性由引理~\ref{lem:2.6.2} 证明).当 $\sum_{j=1}^{\infty}(1 - |a_{j}|) < \infty$ 时,
\begin{equation*}
B(z) = z^{m}\prod_{j = 1}^{\infty}\frac{|a_{j}|}{a_{j}}\frac{a_{j} - z}{1 - \bar{a}_{j}z}
\end{equation*}
称为 Blaschke 积.这里 $m$ 是非负整数,且允许 $a_{j} = a_{i}$($j \neq i$).

当 $0 < x < 1$ 时 $\lim_{x \to 1}(1 - x)/(-\log x) = 1$,故 Blaschke 条件 $\sum_{j=1}^{\infty}(1 - |a_{j}|) < \infty$ 与 $\sum_{j=1}^{\infty}\log 1/|a_{j}| < \infty$ 等价.

\begin{theorem}\label{thm:2.6.1}
对 $\{a_{j}\}_{j=1}^{\infty}$,设 $B$ 为其 Blaschke 积,则 $B \in H^{\infty}$,$|B(z)| < 1$($z \in D$),且 $B$ 的径向极限满足 $|B(e^{i\theta})| = 1$ a.e. $\theta$.$\{a_{j}\}_{j=1}^{\infty}$ 是 $B$ 的零点,$B$ 在 $D \setminus \{0\}$ 上无其他零点.
\end{theorem}

\begin{proof}
若 $\sum_{j=1}^{\infty}(1 - |a_{j}|) < \infty$,则如上所述,无穷乘积 $B$ 在 $D$ 的任意紧集上一致收敛,故 $B \in H(D)$.此时 $B \in H^{\infty}$ 且 $|B(z)| < 1$($z \in D$)是显然的.又可以把 $\{a_{j}\}_{j=1}^{\infty} \cup \{0\}$ 视为 $B$ 的零点集,故设 $B(0) \neq 0$,则假设 $\sum_{j=1}^{\infty}(1 - |a_{j}|) < \infty$ 蕴含 $\log|B(0)| = \sum_{j=1}^{\infty}\log|a_{j}| > -\infty$.此时,由 Jensen 公式(命题~\ref{prop:2.2.1}),
\begin{align*}
\frac{1}{2\pi}\int_{-\pi}^{\pi}\log|B(re^{i\theta})|\dd\theta & = \sum_{|a_{j}| < r}\log\left(\frac{r}{|a_{j}|}\right) + \log|B(0)| \\
& = \sum_{|a_{j}| < r}\log\left(\frac{r}{|a_{j}|}\right) - \sum_{j = 1}^{\infty}\log\frac{1}{|a_{j}|}.
\end{align*}
对任意 $\varepsilon > 0$,存在 $n$ 使得 $\sum_{j > n}\log 1/|a_{j}| < \varepsilon$.把 $r$ 取得足够接近 1 使 $|a_{j}| < r$($j = 1, 2, \ldots, n$),则
\begin{equation*}
\frac{1}{2\pi}\int_{-\pi}^{\pi}\log|B(re^{i\theta})|\dd\theta \geq \sum_{j = 1}^{n}\log\frac{r}{|a_{j}|} - \sum_{j = 1}^{n}\log\frac{1}{|a_{j}|} - \varepsilon.
\end{equation*}
这里,再把 $r$ 取得足够接近 1,则
\begin{equation*}
\frac{1}{2\pi}\int_{-\pi}^{\pi}\log|B(re^{i\theta})|\dd\theta \geq -2\varepsilon,
\end{equation*}
故令 $\varepsilon \to 0$,
\begin{equation*}
\varlimsup_{r \to 1}\frac{1}{2\pi}\int_{-\pi}^{\pi}\log|B(re^{i\theta})|\dd\theta \geq 0.
\end{equation*}
由引理~\ref{lem:2.5.2},$\lim_{r \to 1}B(re^{i\theta}) = B(e^{i\theta})$ a.e. $\theta$ 存在,且 $\log|B(e^{i\theta})| \leq 0$ a.e. $\theta$.由 Lebesgue 积分的 Fatou 引理,
\begin{align*}
0 &\leq \frac{1}{2\pi}\int_{-\pi}^{\pi}\bigl(-\log|B(e^{i\theta})|\bigr)\dd\theta \\
&\leq \varliminf_{r \to 1}\frac{1}{2\pi}\int_{-\pi}^{\pi}\bigl(-\log|B(re^{i\theta})|\bigr)\dd\theta \\
&= -\varlimsup_{r \to 1}\frac{1}{2\pi}\int_{-\pi}^{\pi}\log|B(re^{i\theta})|\dd\theta \leq 0
\end{align*}
故 $|B(e^{i\theta})| = 1$ a.e. $\theta$.
\end{proof}

若 $\{a_{j}\}_{j=1}^{\infty} \subset D$ 满足 $\sum_{j=1}^{\infty}(1 - |a_{j}|) < \infty$,则定理~\ref{thm:2.6.1} 表明存在 $f \in H^{\infty}$ 以其零点集为 $\{a_{j}\}_{j=1}^{\infty}$.其逆是否成立是自然的问题,下面回答这个问题.

\begin{lemma}\label{lem:2.6.2}
设 $\{a_{j}\}_{j=1}^{\infty} \subset D$ 为 $f \in H(D)$ 的零点集.若 $f(z) \not\equiv 0$ 且
\begin{equation*}
\sup_{0 \leq r < 1}\int_{-\pi}^{\pi}\log^{+}|f(re^{i\theta})|\dd\theta < \infty,
\end{equation*}
则 $\sum_{j=1}^{\infty}(1 - |a_{j}|) < \infty$.
\end{lemma}

\begin{proof}
若 $f(0) = 0$,则对 $f(z)/z^{\ell}$ 证明即可,故设 $f(0) \neq 0$.若 $0 < r < 1$ 且 $|a_{j}| \neq r$($j = 1, 2, \ldots$),则由 Jensen 公式(命题~\ref{prop:2.2.1}),
\begin{equation*}
\log|f(0)| = \frac{1}{2\pi}\int_{-\pi}^{\pi}\log|f(re^{i\theta})|\dd\theta + \sum_{|a_{j}| < r}\log\left|\frac{a_{j}}{r}\right|.
\end{equation*}
令 $M = \sup_{0 \leq r < 1}\frac{1}{2\pi}\int_{-\pi}^{\pi}\log^{+}|f(re^{i\theta})|\dd\theta$,则
\begin{equation*}
\sum_{|a_{j}| < r}\log\left|\frac{r}{a_{j}}\right| \leq M - \log|f(0)|,
\end{equation*}
令 $r \to 1$,对任意 $k$,
\begin{equation*}
\sum_{j = 1}^{k}\log\frac{1}{|a_{j}|} \leq M - \log|f(0)|,
\end{equation*}
故
\begin{equation*}
\sum_{j = 1}^{\infty}\log\frac{1}{|a_{j}|} < \infty.
\end{equation*}
\end{proof}

\begin{theorem}\label{thm:2.6.3}
设 $0 < p \leq \infty$.若 $\{a_{j}\}_{j=1}^{\infty}$ 为 $f \in H^{p}$ 的零点集,则 $\sum_{j=1}^{\infty}(1 - |a_{j}|) < \infty$.反之,若 $\sum_{j=1}^{\infty}(1 - |a_{j}|) < \infty$,则存在以 $\{a_{j}\}_{j=1}^{\infty}$ 为零点集的 $f \in H^{p}$.
\end{theorem}

\begin{proof}
可以设 $f(0) \neq 0$.由定理~\ref{thm:2.5.3} 的证明,若 $f \in H^{p}$,则 $f$ 满足引理~\ref{lem:2.6.2} 的假设,故 $\sum_{j=1}^{\infty}(1 - |a_{j}|) < \infty$.反之,若 $\sum_{j=1}^{\infty}(1 - |a_{j}|) < \infty$,由定理~\ref{thm:2.6.1},取 $\{a_{j}\}_{j=1}^{\infty}$ 的 Blaschke 积 $B$ 作为 $f$ 即可.
\end{proof}

\begin{theorem}\label{thm:2.6.4}
设 $0 < p \leq \infty$.若 $f \in H^{p}$ 且 $f(z) \not\equiv 0$,则 $f$ 可以分解为 $f(z) = B(z)g(z)$.这里 $B$ 是 Blaschke 积,$g$ 在 $D$ 内无零点且属于 $H^{p}$.
\end{theorem}

\begin{proof}
可以设 $f(0) \neq 0$ 且其零点集为无限集.设 $\{a_{j}\}_{j=1}^{\infty}$ 为 $f \in H^{p}$ 的零点集,由定理~\ref{thm:2.6.3} 得 $\sum_{j=1}^{\infty}(1 - |a_{j}|) < \infty$.由定理~\ref{thm:2.6.1} 存在 $\{a_{j}\}_{j=1}^{\infty}$ 的 Blaschke 积 $B$,令 $g(z) = f(z)/B(z)$,则 $g$ 在 $D$ 内无零点.还需证明 $g \in H^{p}$.

令
\begin{equation*}
B_{n}(z) = \prod_{j = 1}^{n}\frac{|a_{j}|}{a_{j}}\frac{a_{j} - z}{1 - \bar{a}_{j}z} \quad \text{且} \quad g_{n}(z) = f(z)/B_{n}(z),
\end{equation*}
则对每个 $n$ 和任意 $\varepsilon > 0$,当 $|z|$ 足够接近 1 时可使 $|B_{n}(z)| > 1 - \varepsilon$.因此,当 $r$ 足够接近 1 时,
\begin{align*}
\frac{1}{2\pi}\int_{-\pi}^{\pi}|g_{n}(re^{i\theta})|^{p}\dd\theta & \leq (1 - \varepsilon)^{-p}\frac{1}{2\pi}\int_{-\pi}^{\pi}|f(re^{i\theta})|^{p}\dd\theta \\
& \leq (1 - \varepsilon)^{-p}\sup_{0 \leq r < 1}M_{p}(r, f)^{p}.
\end{align*}
令 $\varepsilon \to 0$,当 $r$ 足够接近 1 时,
\begin{equation*}
\frac{1}{2\pi}\int_{-\pi}^{\pi}|g_{n}(re^{i\theta})|^{p}\dd\theta \leq \sup_{0 \leq r < 1}M_{p}(r, f)^{p}.
\end{equation*}
由引理~\ref{lem:2.1.1},可证当 $r' < r$ 时 $M_{p}(r', g_{n}) \leq M_{p}(r, g_{n})$,故上面的不等式对所有 $r$ 成立.因此由 Fatou 引理,
\begin{equation*}
\sup_{0 \leq r < 1}M_{p}(r, g) \leq \sup_{0 \leq r < 1}M_{p}(r, f) < \infty
\end{equation*}
得证.
\end{proof}
